\documentclass[10pt,a4paper]{amsart}
\usepackage[margin=19mm]{geometry}
\usepackage{amsmath,amssymb,mathtools}
\usepackage[scr=rsfs,cal=euler]{mathalfa}
\usepackage{xfrac}
\usepackage[unicode,hidelinks,bookmarks=false]{hyperref}
\newcommand{\ie}{i.e.\ }
\newcommand{\cf}{cf.\ }
\newcommand{\resp}{respectively\ }
\newcommand{\Subsection}{\subsection}
\providecommand{\nobreakdash}{\nobreak\hbox{-}}
\setlength{\emergencystretch}{2em}
\multlinegap0pt
\usepackage{amssymb}
\usepackage{mathtools}
\usepackage{algorithm}\usepackage{algpseudocode}
\usepackage{xcolor}
\newtheorem{theorem}{Theorem}
\title{A Control-Theoretic Approach for Resource-Aware Consensus in Multi-Agent AI}
\author{James Flagg, Esteban A. Hernandez-Vargas}
\date{Source: arXiv:2608.25099v1 (CC BY 4.0)}
\begin{document}
\maketitle

\noindent\emph{Source and licence.} A Control-Theoretic Approach for Resource-Aware Consensus in Multi-Agent AI. Version arXiv:2608.25099v1 (CC BY 4.0), distributed by arXiv under the Creative Commons Attribution 4.0 International licence, http://creativecommons.org/licenses/by/4.0/, as recorded in the arXiv licence metadata for this exact version and shown on the abstract page https://arxiv.org/abs/2608.25099v1. Full licence notice: \texttt{licenses/arxiv-2608.25099-CC-BY-4.0.txt}.\par\smallskip

\noindent\textit{Editorial scope.} Counted unit: the article's theorem thm:rac (source0.tex lines 606--671). The article's complete original proof of it is delivered with it (source0.tex lines 673--691, one \texttt{proof} environment of 19 lines). No mathematics is written, repaired or re-derived: the statement and the proof are the article's own lines, in the article's own notation, and no retained line is substituted or re-flowed.  Retained uncounted as the source-local context the proof needs: lines 460--464.  Nothing was omitted from any retained span, so the package carries no omission record.  No retained line contains a citation, so the wrapper carries no bibliography and no bibliography entry.  No retained line contains a figure environment, an image-inclusion command or drawing code, so no image is delivered and the package declares no asset dependency.  Editorial formatting only, in addition: an AMS \texttt{amsart} preamble that restates the article's own declarations and macros used by the retained lines, namely the environments \texttt{theorem} on separate counters, together with the standard packages those lines need.  The retained environments print in this document as Theorem 1 in order of appearance, whereas the article prints the same items with the numbering of its own full document.  The complete native source of the article as distributed in the arXiv e-print for this exact version is bundled unchanged as \texttt{sources/208547/source0.tex}.
\begin{equation} \label{eq:contracting}
D(k+1)
\le
\lambda_{\sigma_k}D(k),
\end{equation}

\begin{theorem}[Resource-Aware Consensus]
\label{thm:rac}

Assume $ D(0)>\eta>\varepsilon>0.$ Define
\begin{equation}
K_1
:=
\left\lceil
\frac{\log(\eta/D(0))}
{\log\lambda_1}
\right\rceil,
\qquad
K_2
:=
\left\lceil
\frac{\log(\varepsilon/\eta)}
{\log\lambda_2}
\right\rceil,
\label{eq:Ks}
\end{equation}
for consensus-contraction rates $0< \lambda_1, \lambda_2 < 1 $,
and let $K^* = K_1+K_2$.


Then:

\begin{enumerate}
\item
During the high-disagreement phase,
\[
D(k)
\le
\lambda_1^kD(0),
\qquad
0< k < K_1.
\]

\item
During the low-disagreement phase,
\[
D(k)
\le
\lambda_2^{\,k-K_1}\eta,
\qquad
k\geq K_1>0.
\]

\item
The trajectory reaches
$\varepsilon$-consensus in at most $K^*$ rounds:
\[
D(K^*)
\le
\varepsilon.
\]

\item
The corresponding communication cost is
\[
B^*
=
K_1c_1+K_2c_2,
\] where $c_1,c_2$ are the per-round token costs of modes 1 and 2 respectively.
\end{enumerate}

\end{theorem}

\begin{proof}
To derive the expression for $K_1$, observe that repeated application of \eqref{eq:contracting} yields $D(K_1) \leq \lambda_1^{K_1}D(0)$ for $K_1 > 0.$
The switching rule activates mode $1$ while
$D(k)>\eta$ and mode $2$ otherwise.
$K_1$ is an upper bound on the number of rounds until $D(k) \leq \eta$. To solve for $K_1$, set $\eta = \lambda_1^{K_1}D(0) \implies \lambda_1^{K_1} = \frac{\eta}{D(0)} \implies K_1 = \log_{\lambda_1}\Big(\frac{\eta}{D(0)} \Big)$. Changing base, $K_1 = \frac{\log (\eta / D(0))}{\log \lambda_1}$. Since we are counting discrete communication rounds, we must have $K_1 \in \mathbb{Z}_{\geq 0}$. We also want to guarantee $D(K_1) \leq \eta $. Thus, we apply the ceiling function $K_1
=
\left\lceil
\frac{\log(\eta/D(0))}
{\log\lambda_1}
\right\rceil.$
The expression for $K_2$ is derived similarly. Due to the construction of $K_1,K_2$, repeated application of \eqref{eq:contracting} yields the bounds in (1) and (2).  The definition of $K_1$ implies
$ \lambda_1^{K_1}D(0)\le\eta$, while the definition of $K_2$ implies $ \lambda_2^{K_2}\eta\le\varepsilon$.
Combining both estimates yields $D(K^*)\le\varepsilon$.

Finally, the total token expenditure equals the number of rounds
spent in each mode multiplied by the corresponding communication cost,
giving $B^*=K_1c_1+K_2c_2$.

\end{proof}

\end{document}
